% !TeX root = main.tex

% Аргументы: имя файла в images, подпись, метка для ссылки.
% После добавления файла рамка автоматически заменяется изображением.
\newcommand{\imageplaceholder}[3]{%
  \begin{figure}[!htbp]
    \centering
    \IfFileExists{images/#1}{%
      \includegraphics[width=0.95\textwidth,keepaspectratio]{#1}%
    }{%
      \fbox{%
        \parbox[c][3cm][c]{\dimexpr0.95\textwidth-2\fboxsep-2\fboxrule\relax}{%
          \centering
          Место для рисунка\\[0.4em]
          \texttt{images/\detokenize{#1}}
        }%
      }%
    }
    \caption{#2}
    \label{#3}
  \end{figure}
}

% Видимые места для текста, который будет добавлен после выполнения работы.
\newcommand{\reportplaceholder}[1]{%
  \par\noindent\textit{[#1]}\par
}
\newcommand{\fillvalue}{\textit{указать}}

\tableofcontents
\newpage

\section{Цель работы}

Освоить метод расширенной ошибки в задаче адаптивного управления
линейным объектом по выходу. Исследовать стабилизацию и слежение
при различных коэффициентах адаптации, а также работу регулятора
после изменения знаков коэффициентов объекта.

\section{Постановка задачи и исходные данные}

Рассматривается линейный объект второго порядка
\begin{equation}
  \ddot y+a_1\dot y+a_0y=b_0u,
  \qquad n=2,\qquad m=0.
  \label{eq:plant}
\end{equation}
Регулятору доступны выход $y(t)$ и формируемое управление $u(t)$.
Коэффициенты $a_0,a_1,b_0$ считаются неизвестными регулятору;
их истинные значения используются при моделировании объекта.
Предполагается, что известна нижняя граница
$b_0\ge b_{\min}>0$.

Исходные данные:
\begin{equation}
  a_0=1,\qquad a_1=0,\qquad b_0=1,
  \qquad k_{M1}=2,\qquad k_{M0}=1,\qquad k_0=1.
  \label{eq:variant}
\end{equation}
Следовательно, уравнение и передаточная функция объекта имеют вид
\begin{equation}
  \ddot y+y=u,
  \qquad W(s)=\frac{1}{s^2+1}.
  \label{eq:variant-plant}
\end{equation}
Полиномы эталонной модели и вспомогательных фильтров:
\begin{equation}
  K_M(s)=s^2+2s+1=(s+1)^2,
  \qquad K(s)=s+1.
  \label{eq:polynomials}
\end{equation}
Задающее воздействие в опытах слежения:
\begin{equation}
  g(t)=2\sgn(\cos t).
  \label{eq:reference-input}
\end{equation}
При исследовании стабилизации принимается $g(t)\equiv0$.

Ошибка слежения и цель управления определяются соотношениями
\begin{equation}
  \varepsilon(t)=y_M(t)-y(t),
  \qquad \lim_{t\to\infty}\varepsilon(t)=0,
  \label{eq:goal}
\end{equation}
при ограниченности всех сигналов замкнутой системы.

\section{Построение адаптивной системы управления}

\subsection{Эталонная модель и вспомогательные фильтры}

Эталонная модель согласно (7.5) методических указаний:
\begin{equation}
  y_M=\frac{k_0}{K_M(s)}[g]
  =\frac{1}{(s+1)^2}[g],
  \qquad \ddot y_M+2\dot y_M+y_M=g.
  \label{eq:reference-model}
\end{equation}
Фильтры (7.2), (7.3) для выбранного полинома $K(s)=s+1$:
\begin{equation}
  \dot v_1=-v_1+u,
  \qquad \dot v_2=-v_2+y,
  \qquad
  \bm\omega=\begin{pmatrix}v_1&v_2&y\end{pmatrix}^{\mathsf T}.
  \label{eq:filters}
\end{equation}

\reportplaceholder{Привести параметризованное представление объекта
по формуле (7.4) и определить компоненты вектора $\bm\psi$.}

\subsection{Настраиваемый регулятор}

Закон управления (7.8) имеет вид
\begin{equation}
  u=\frac{-\hat{\bm\psi}^{\mathsf T}\bm\omega+g}{\hat b_0},
  \qquad
  \hat{\bm\psi}_p=
  \begin{pmatrix}
    \hat\psi_1&\hat\psi_2&\hat\psi_3&\hat b_0
  \end{pmatrix}^{\mathsf T}.
  \label{eq:control}
\end{equation}
Для предотвращения деления на ноль используется ограничение
$\hat b_0(t)\ge b_{\min}$.

\reportplaceholder{Описать используемую проекцию оценки $\hat b_0$
и указать выбранное значение $b_{\min}$.}

\subsection{Расширенная ошибка и алгоритм адаптации}

\reportplaceholder{Ввести регрессор $\bm\omega_p$, фильтрованные сигналы
и расширенную ошибку $\hat\varepsilon$ по формуле (8.2).
Привести реализацию дополнительных фильтров.}

\reportplaceholder{Записать принятый алгоритм адаптации для
$\hat{\bm\psi}$ и $\hat b_0$ с коэффициентами $\gamma_1$, $\gamma_2$.
Указать используемый нормирующий множитель.}

% При заполнении сверить нормировку: в (8.4) используется
% 1 + bar(omega_p)^T bar(omega_p), а в (8.5), (8.6) записано
% 1 + bar(omega)^T bar(omega). Это разные выражения.

\reportplaceholder{Пояснить назначение расширенной ошибки
и условия сходимости ошибки слежения.}

\subsection{Условия моделирования}

\begin{table}[!htbp]
  \centering
  \caption{Общие параметры моделирования}
  \label{tab:simulation-settings}
  \begin{tabularx}{\textwidth}{@{}Xl@{}}
    \toprule
    Параметр & Значение \\
    \midrule
    Продолжительность моделирования $T$, с & \fillvalue \\
    Решатель и настройки точности & \fillvalue \\
    Нижняя граница $b_{\min}$, $0<b_{\min}\le1$ & \fillvalue \\
    Начальные $y(0)$, $\dot y(0)$ & \fillvalue \\
    Начальные $y_M(0)$, $\dot y_M(0)$ & \fillvalue \\
    Начальные $v_1(0)$, $v_2(0)$ & \fillvalue \\
    Начальная оценка $\hat{\bm\psi}(0)$ & \fillvalue \\
    Начальная оценка $\hat b_0(0)$ & \fillvalue \\
    Начальные состояния дополнительных фильтров & \fillvalue \\
    \bottomrule
  \end{tabularx}
\end{table}

По заданию в опыте стабилизации выбираются ненулевые начальные
условия фильтров~\eqref{eq:filters} и алгоритма настройки
$\hat{\bm\psi}$, а начальная оценка коэффициента управления
удовлетворяет $\hat b_0(0)>b_0=1$.
В каждом из первых двух опытов рассматриваются три различные пары
коэффициентов; $\gamma_1$ выбирается на несколько порядков больше
$\gamma_2$.

\reportplaceholder{Уточнить начальные условия для каждого опыта,
если они отличаются от приведённых в таблице~\ref{tab:simulation-settings}.}

\newpage
\section{Исследование стабилизирующего адаптивного управления}

В первом опыте принимается $g(t)\equiv0$.
Исследуется стремление выхода объекта к нулю при трёх парах
коэффициентов адаптации.

\begin{table}[!htbp]
  \centering
  \caption{Коэффициенты адаптации в опыте стабилизации}
  \label{tab:stabilization-gains}
  \begin{tabular}{@{}ccc@{}}
    \toprule
    Номер пары & $\gamma_1$ & $\gamma_2$ \\
    \midrule
    1 & \fillvalue & \fillvalue \\
    2 & \fillvalue & \fillvalue \\
    3 & \fillvalue & \fillvalue \\
    \bottomrule
  \end{tabular}
\end{table}

\subsection{Результаты моделирования}

Для сравнения трёх пар коэффициентов предусмотрены графики
выхода $y$, управления $u$ и четырёх оценок
$\hat{\bm\psi}_p$ на
рисунках~\ref{fig:stabilization-y}--\ref{fig:stabilization-parameters}.

\imageplaceholder{stabilization_y.png}
  {Выход объекта при стабилизации для трёх пар коэффициентов адаптации}
  {fig:stabilization-y}

\imageplaceholder{stabilization_u.png}
  {Управляющее воздействие при стабилизации для трёх пар коэффициентов адаптации}
  {fig:stabilization-u}

\imageplaceholder{stabilization_parameters.png}
  {Оценки параметров при стабилизации для трёх пар коэффициентов адаптации}
  {fig:stabilization-parameters}

\newpage
\subsection{Анализ результатов}

\reportplaceholder{Сравнить переходные процессы при трёх парах
коэффициентов: затухание выхода, величину и характер управления,
поведение оценок параметров. Выводы сформулировать по полученным графикам.}

\noindent\textbf{Вывод по пункту:}
\reportplaceholder{Сформулировать результат опыта стабилизации.}

\section{Исследование следящего адаптивного управления}

Во втором опыте используется задание~\eqref{eq:reference-input}.
Выход объекта сравнивается с выходом эталонной
модели~\eqref{eq:reference-model} при трёх парах коэффициентов адаптации.

\begin{table}[!htbp]
  \centering
  \caption{Коэффициенты адаптации в опыте слежения}
  \label{tab:tracking-gains}
  \begin{tabular}{@{}ccc@{}}
    \toprule
    Номер пары & $\gamma_1$ & $\gamma_2$ \\
    \midrule
    1 & \fillvalue & \fillvalue \\
    2 & \fillvalue & \fillvalue \\
    3 & \fillvalue & \fillvalue \\
    \bottomrule
  \end{tabular}
\end{table}

\subsection{Результаты моделирования}

Для сравнения трёх пар коэффициентов предусмотрены графики
$y$ и $y_M$, управления $u$ и оценок $\hat{\bm\psi}_p$
на рисунках~\ref{fig:tracking-y}--\ref{fig:tracking-parameters}.

\imageplaceholder{tracking_y.png}
  {Выходы объекта и эталонной модели при трёх парах коэффициентов адаптации}
  {fig:tracking-y}

\imageplaceholder{tracking_u.png}
  {Управляющее воздействие при слежении для трёх пар коэффициентов адаптации}
  {fig:tracking-u}

\imageplaceholder{tracking_parameters.png}
  {Оценки параметров при слежении для трёх пар коэффициентов адаптации}
  {fig:tracking-parameters}

\newpage
\subsection{Анализ результатов}

\reportplaceholder{Сравнить качество слежения и переходные процессы
при трёх парах коэффициентов. Проанализировать управление и оценки
параметров, различая сходимость ошибки слежения и ошибок параметров.}

\noindent\textbf{Вывод по пункту:}
\reportplaceholder{Сформулировать результат опыта слежения.}

\section{Слежение после изменения знаков коэффициентов объекта}

В третьем опыте знаки $a_0$ и $a_1$ меняются на противоположные.
Для варианта 1 коэффициент $a_1=0$ сохраняется, поэтому
\begin{equation}
  a_0=-1,\qquad a_1=0,\qquad b_0=1,
  \qquad \ddot y-y=u.
  \label{eq:modified-plant}
\end{equation}
Передаточная функция изменённого объекта имеет вид
\begin{equation}
  W_{\text{изм}}(s)=\frac{1}{s^2-1}.
\end{equation}
Сохраняются задание~\eqref{eq:reference-input} и
эталонная модель~\eqref{eq:reference-model}.
Моделирование проводится для одной пары коэффициентов адаптации.

\reportplaceholder{Указать выбранные $\gamma_1$, $\gamma_2$
и начальные условия.}

\subsection{Результаты моделирования}

Для изменённого объекта предусмотрены графики $y$ и $y_M$,
управления $u$ и оценок $\hat{\bm\psi}_p$
на рисунках~\ref{fig:modified-y}--\ref{fig:modified-parameters}.

\imageplaceholder{modified_plant_y.png}
  {Выходы изменённого объекта и эталонной модели}
  {fig:modified-y}

\imageplaceholder{modified_plant_u.png}
  {Управляющее воздействие после изменения знаков коэффициентов объекта}
  {fig:modified-u}

\imageplaceholder{modified_plant_parameters.png}
  {Оценки параметров после изменения знаков коэффициентов объекта}
  {fig:modified-parameters}

\newpage
\subsection{Анализ результатов}

\reportplaceholder{Сопоставить результаты с опытом слежения
для исходного объекта. Оценить влияние изменения знаков коэффициентов
на выход, управление и настройку параметров.}

\noindent\textbf{Вывод по пункту:}
\reportplaceholder{Сформулировать результат опыта с изменённым объектом.}

\section{Выводы}

\reportplaceholder{По фактическим результатам обобщить работу метода
расширенной ошибки, влияние коэффициентов адаптации на стабилизацию
и слежение, а также поведение системы при изменении параметров объекта.
Указать наблюдаемые ограничения и различие между достижением цели
управления и точностью оценивания параметров.}
